\chapter*{Словарь терминов}             % Заголовок
\addcontentsline{toc}{chapter}{Словарь терминов}  % Добавляем его в оглавление

\textbf{RTT} -  round-trip-time, <<время круговой передачи>>. Время от отправки сообщения адресату до получения от него ответа.

\textbf{FLOPs (Floating Point Operations)} - это число операций с плавающей точкой, необходимых для выполнения одного прохода модели машинного обучения. Эта метрика отражает вычислительную сложность и используется для оценки эффективности и ресурсоёмкости алгоритма.

\textbf{Accuracy (Точность)} - метрика качества модели в задачах классификации. Отражает долю правильных предсказаний модели среди всех сделанных предсказаний.

\textbf{Precision (Точность)} - метрика качества модели в задачах классификации. Доля истинно положительных предсказаний среди всех положительных
предсказаний модели.

\textbf{Recall (Полнота)}  - метрика качества модели в задачах классификации. Доля истинно положительных предсказаний среди всех реально положительных случаев.

\textbf{F1-метрика} - метрика качества модели в задачах классификации. Гармоническое среднее между Precision и Recall, учитывает обе метрики.

\textbf{p-value} - это вероятность получить наблюдаемое (или более экстремальное) значение статистики, при условии, что нулевая гипотеза верна.

\textbf{pdf} - (расшифровка Probability Density Function) функция плотности вероятности, которая описывает непрерывное распределение вероятностей случайной величины.